\section{模型假设}

\begin{enumerate}
\item \textbf{损失可加假设}：验证损失 $L$ 由不可约误差 $E$、参数量不足、
数据量不足与质量不足四项可加贡献构成，即广义标度律的可加分解形式成立；
\item \textbf{质量指标独立假设}：各质量指标在同一文档上独立观测，可用加权
求和或 TOPSIS 合成；指标间冲突可用``内容质量''与``格式清洁''两族均值之差的
绝对值显式刻画；
\item \textbf{质量成本可分离假设}：质量提升成本 $C_Q$ 仅与质量水平 $Q$、
基线 $Q_0$ 与数据量 $D$ 有关，且只在 $Q>Q_0$ 时产生，形式为
$C_Q=D[g(Q)-g(Q_0)]_+$；
\item \textbf{外生上下文长度假设}：$L_{\text{ctx}}$ 为模型架构确定的外生量，
不参与优化，其可行取值依据附件 C7 的 \texttt{max\_position\_embeddings}；
\item \textbf{训练成本常数假设}：训练 FLOPs 按 $C_{\text{train}}=6ND$ 估计
（每参数每 token 约 6 FLOPs），注意力成本按系数 $\eta=2\times10^{-4}$ 计入；
\item \textbf{损失可比假设}：同一规模族内各验证损失在相同验证集上可比；
跨家族比较时需先做家族偏移修正；
\item \textbf{配比自由假设}：调整领域配比 $p$ 本身不额外消耗算力，配比效应
通过加权平均质量 $Q=\sum_i p_i q_i$ 进入能力方程；
\item \textbf{前沿外推假设}：前沿能力的对数与参数对数、时间近似线性，该关系
在未来 24 个月内保持（不确定性通过自助区间刻画）。
\end{enumerate}
